\documentclass[../main.tex]{subfiles}
\begin{document}

Chương 1 đã chỉ ra rằng các kiến trúc agent hiện có cho \gls{ISD} chủ yếu tối ưu cho chất lượng artifact tương thích ADDIE tại lần sinh đầu tiên, chưa đặt mức độ alignment giữa các thành phần thiết kế làm trạng thái trung tâm có thể truy vết, kiểm tra, sửa chữa hay tinh chỉnh hiệu quả. Chương 2 này khảo sát các công trình liên quan để làm rõ khoảng trống đó từ góc độ kỹ thuật. Vì luận văn được trình bày trước hội đồng Công nghệ thông tin, trọng tâm khảo sát đặt vào kiến trúc agent và benchmark cho LLM; lý thuyết ISD chỉ được tóm lược ở mức tối thiểu cần thiết để đọc phần phương pháp ở Chương 3. Phần \ref{sec:scope} xác định phạm vi nghiên cứu, phần \ref{sec:related-work} khảo sát ba nhóm công trình liên quan theo thứ tự từ kiến trúc agent tổng quát đến các hệ thống ISD cụ thể, và phần cuối tóm lược các khái niệm ISD nền tảng.

\section{Scope of Research}
\label{sec:scope}

Luận văn tiếp cận instructional design như một bài toán kỹ thuật: sinh và duy trì một tập artifact có cấu trúc (mục tiêu học tập, bài đánh giá, hoạt động, tài liệu, kế hoạch triển khai, kế hoạch đánh giá) sao cho các artifact này nhất quán với nhau, dưới sự điều phối của một hệ agent dùng LLM. Phạm vi khảo sát, vì vậy, không mở rộng sang giáo dục học nói chung, mà giới hạn ở ba nhóm công trình có liên quan trực tiếp đến việc thiết kế kiến trúc agent: (i) kiến trúc agent và multi-agent orchestration dùng LLM, (ii) phương pháp benchmark và đánh giá agent, và (iii) các hệ thống agent áp dụng cụ thể cho ISD. Lý thuyết ISD như ADDIE hay constructive alignment chỉ được khảo sát ở mức khái niệm, đóng vai trò ràng buộc miền (domain constraint) mà kiến trúc agent phải thoả mãn, không phải đối tượng nghiên cứu chính của luận văn.

\section{Related Work}
\label{sec:related-work}

\subsection{LLM Agent Architectures and Multi-agent Orchestration}

Trước khi so sánh các chiến lược orchestration cụ thể, cần làm rõ khái niệm harness được dùng xuyên suốt luận văn. Trong các hệ agent dùng LLM, harness là một lớp vỏ nhiều tầng bao quanh một model nền tảng vốn không có trạng thái (stateless), đảm nhiệm việc dựng prompt, quản lý ngữ cảnh và bộ nhớ, gọi công cụ, xử lý lỗi, và áp đặt các ràng buộc an toàn/kiểm tra, thay vì chỉ là một wrapper gọi \gls{API} đơn giản. Vai trò của lớp này ngày càng quan trọng khi năng lực model nền bão hoà trên nhiều benchmark \cite{deveci2025ouroboros,shukla2025scalingwall}: với cùng một model, thay đổi harness có thể dịch chuyển đáng kể kết quả trên các benchmark agent, và tác động của harness lớn nhất chính ở các model nhỏ \cite{cho2026harness}. Hình~\ref{fig:harness-anatomy} minh hoạ khái niệm này ở mức tổng quát: một yêu cầu đi vào runtime (orchestration loop, prompt construction, output parsing, error handling), được bao bọc bởi lớp capability (tool, memory, context/state management) và lớp safety/scale (guardrail, verification loop, sub-agent orchestration).

\begin{figure}[H]
\centering
\includegraphics[width=0.85\linewidth]{Figure/harness-anatomy.jpeg}
\caption{Minh hoạ khái niệm harness nhiều tầng bao quanh một LLM stateless, đối lập với cách hiểu harness chỉ là một wrapper gọi API đơn giản.}
\label{fig:harness-anatomy}
\end{figure}

Các chiến lược orchestration khảo sát dưới đây, single-prompt, ReAct tool-use và multi-agent generate-evaluate-optimize, chính là những cách khác nhau để thiết kế lớp runtime và capability bên trong một harness. Nhóm công trình đầu tiên và có liên quan kỹ thuật gần nhất là các kiến trúc agent dùng LLM để thực hiện tác vụ nhiều bước. Chiến lược đơn giản nhất là single-prompt generation, trong đó toàn bộ output được sinh trong một lượt gọi model; chiến lược này dễ triển khai và chi phí thấp, nhưng không có cơ chế kiểm tra từng phần và khó sửa lỗi cục bộ khi một phần output sai. Chiến lược ReAct xen kẽ bước suy luận (reasoning) và bước hành động (tool-use), cho phép agent gọi công cụ bên ngoài và điều chỉnh kế hoạch dựa trên kết quả quan sát được; trong lĩnh vực ISD, chiến lược này được hiện thực hoá bởi agent \texttt{react-isd} (kiến trúc ReAct-ADDIE) \cite{yao2023react,isdagentbench}. Một hướng khác là multi-agent orchestration, trong đó nhiều agent chuyên biệt theo vai trò (sinh nội dung, đánh giá, tối ưu) phối hợp qua một vòng lặp generate-evaluate-optimize, như trong EduPlanner \cite{zhang2025eduplanner}; cách tiếp cận này cải thiện chất lượng output qua phản hồi lặp lại nhưng thường chưa chuẩn hoá cơ chế truy xuất tri thức theo tác vụ (agentic retrieval) và chưa có schema contract hay trace thống nhất ở mọi bước. Ba chiến lược này, single-prompt, ReAct tool-use và multi-agent generate-evaluate-optimize, đại diện cho phổ kiến trúc agent phổ biến hiện nay và là điểm tham chiếu kỹ thuật chính khi thiết kế AlignmentGraph-ISD ở Chương 3.

\subsection{Benchmarking and Evaluation of LLM Agents}

Nhóm công trình thứ hai tập trung vào phương pháp đánh giá năng lực của LLM và agent trong các tác vụ có yếu tố giáo dục, tiêu biểu là OpenLearnLM \cite{openlearnlm2026}, EduBench \cite{edubench2025}, MathTutorBench \cite{mathtutorbench2025} và TutorBench \cite{tutorbench2025}. Nhìn từ góc độ thiết kế benchmark, các công trình này đóng góp phương pháp luận về cách xây dựng bộ scenario, cách dùng LLM-as-judge để chấm điểm output tự do dạng văn bản, và cách phân tích kết quả theo nhiều trục như độ khó hay lĩnh vực kiến thức. Tuy nhiên, các benchmark này đánh giá năng lực model ở cấp tác vụ đơn lẻ, như trả lời câu hỏi, chấm bài hoặc tutoring hội thoại, mà không đánh giá một quy trình nhiều bước có trạng thái, nơi các artifact được sinh ra phải nhất quán với nhau. ISD-Agent-Bench \cite{isdagentbench}, benchmark được dùng để đánh giá luận văn này (chi tiết ở Chương 5), khác biệt ở chỗ đánh giá đồng thời chất lượng artifact theo rubric ADDIE và chất lượng trajectory của cả quá trình agent thực thi, phù hợp hơn với bài toán ISD nhiều bước mà luận văn đặt ra.

\subsection{ISD-specific Agent Systems}

Nhóm công trình thứ ba, gần với luận văn nhất, là các hệ thống agent áp dụng trực tiếp cho instructional design. LearnLM \cite{learnlm2024} và LessonPlanner \cite{lessonplanner2024} tối ưu hoá model hoặc pipeline cho tác vụ giảng dạy và lập kế hoạch bài học; các nghiên cứu sinh giáo án bằng GPT-4 \cite{hu2024teachingplan} cho thấy LLM có thể tạo giáo án khả dụng nhưng thường dừng ở một artifact đơn lẻ, không đặt trong một vòng đời ISD có kiểm tra xuyên suốt. ISD-Agent-Bench đóng góp thêm bốn agent baseline dựa trên lý thuyết ISD cụ thể: ADDIE-Agent \cite{isdagentbench,branch2009addie} hiện thực một StateGraph tuyến tính qua năm pha ADDIE; Dick-Carey-Agent \cite{isdagentbench,dick1978systematic} theo systems approach của Dick và Carey với vòng lặp formative-evaluation; RPISD-Agent \cite{isdagentbench,tripp1990rapid} theo rapid prototyping với phản hồi nhiều phía; và ReAct-ADDIE/\texttt{react-isd} \cite{isdagentbench,yao2023react} như đã nêu ở trên. Bảng~\ref{tab:agent-comparison} tóm tắt kiến trúc lõi, điểm mạnh, giới hạn chính của từng hệ thống, và khác biệt tương ứng của kiến trúc đề xuất trong luận văn.

\begin{table}[H]
\centering
\small
\caption{So sánh kiến trúc AlignmentGraph-ISD với các agent trong ISD-Agent-Bench.}
\label{tab:agent-comparison}
\begin{tabular}{p{2.1cm}p{2.6cm}p{2.5cm}p{2.8cm}p{3.2cm}}
\hline
\textbf{Agent} & \textbf{Kiến trúc lõi} & \textbf{Điểm mạnh} & \textbf{Giới hạn chính} & \textbf{Khác biệt của AlignmentGraph-ISD} \\
\hline
\texttt{baseline} & Một prompt sinh toàn bộ ADDIE output & Đơn giản, chi phí thấp & Không có kiểm tra pha, khó sửa lỗi cục bộ & Chia workflow thành node/edge, validate từng phần, log trajectory \\
\texttt{react-isd} & ReAct agent dùng ISD tools & Linh hoạt, có tool-use & Thứ tự và chất lượng phụ thuộc mạnh vào model & Dùng graph policy kiểm soát transition và tool permission \\
\texttt{addie-agent} & StateGraph tuyến tính 5 pha ADDIE & Rõ lifecycle, dễ tương thích schema & Tuyến tính cứng, ít cơ chế repair & Giữ ADDIE contract nhưng cho phép dynamic orchestration và repair \\
\texttt{dick-carey-agent} & Systems approach, có formative-evaluation loop & Mạnh ở revision theo hệ thống & Gắn cứng một process model & Không gắn cứng một model ISD, dùng policy/verifier kết hợp \\
\texttt{rpisd-agent} & Rapid prototyping, phản hồi nhiều phía & Phù hợp cải tiến prototype & Phức tạp, không phải scenario nào cũng cần & Bật/tắt repair loop theo verifier result \\
\texttt{eduplanner} & Multi-agent generate-evaluate-optimize & Có evaluator/optimizer loop & Chưa chuẩn hoá agentic retrieval và trace & Multi-agent theo vai trò ISD, context pack theo task, trace thống nhất \\
\hline
\end{tabular}
\end{table}

\section{Background Concepts}

Vì phạm vi nghiên cứu là kỹ thuật, phần này chỉ tóm lược ngắn gọn hai khái niệm ISD cần thiết để đọc Chương 3, không mở rộng thảo luận sư phạm. ADDIE \cite{branch2009addie} là mô hình quy trình instructional design gồm năm pha Analysis, Design, Development, Implementation và Evaluation; luận văn dùng ADDIE làm output contract để tương thích với ISD-Agent-Bench, không dùng làm ràng buộc thứ tự thực thi của agent. Constructive alignment \cite{biggs1996enhancing} là nguyên lý cho rằng learning objectives, learning activities và assessment tasks phải nhất quán với nhau; đây là cơ sở lý thuyết trực tiếp cho hai quan hệ \texttt{measures} và \texttt{supports} trong alignment graph ở Chương 3. Systems approach của Dick và Carey \cite{dick1978systematic} bổ sung một vòng formative-evaluation/revision cho quy trình thiết kế; luận văn dùng nguyên lý này để biện minh cho cơ chế cascade-invalidation, thay vì hiện thực hoá toàn bộ quy trình Dick-Carey.

Ba khái niệm trên, ADDIE làm output contract, constructive alignment làm cơ sở cho alignment edge, và systems approach làm cơ sở cho cascade-invalidation, là toàn bộ nền tảng lý thuyết ISD mà Chương 3 sử dụng. Kết hợp với ba nhóm công trình kỹ thuật đã khảo sát ở trên, có thể thấy khoảng trống chung: các kiến trúc agent hiện có hoặc không kiểm tra alignment một cách tường minh, hoặc gắn cứng vào một model ISD cụ thể, hoặc chưa chuẩn hoá cơ chế retrieval/schema/trace xuyên suốt toàn bộ vòng đời thiết kế; và hầu hết đều tối ưu cho lần sinh đầu tiên thay vì vòng tinh chỉnh nhiều lượt. Chương 3 trình bày cách AlignmentGraph-ISD giải quyết các hạn chế này bằng một cấu trúc dữ liệu graph tường minh và một harness thực thi phi tuyến xung quanh cấu trúc đó.

\end{document}
